% !TeX program = xelatex
% Main file: long_report.tex. Style preserved from the assignment template.
\documentclass[11pt]{article}
\XeTeXinputnormalization=1
\usepackage[ko]{mas1004}
% Keep PDF navigation available without hyperref's extra bookmark pass.
\usepackage{bookmark}
% Use fast lossless PDF stream compression; image pixels are unchanged.
\AddToHook{shipout/firstpage}{\special{dvipdfmx:config z 1}}
\setmonohangulfont{Noto Sans Mono CJK KR}[Script=Hangul,Language=Korean]
\reporttitle{Assignment 1: Long report}
\studentname{김서영}
\studentid{20251120}
\begin{document}
\makeheader

\section*{Problem 1. Categories and images}
\subsection*{Why these categories}
분류 대상은 명조의 양양, 카멜리아, 린네, 에이메스이다.
루파는 보유 사진이 적어 제외했고, 페비 대신 린네로 변경했다.
공식 일러스트, 팬 일러스트, 인게임 스크린샷을 모두 포함한다.
선정 이유는 학생이 다음과 같이 설명했다.
\begin{quote}
``내가 평소에 매우 좋아하는 명조와 관련된 것을 일단 하고 싶었다. 그리고 명조 일러스트가 다양한 그림체, 3D 모델이 많기에 그것들을 분류하고 싶었다.''
\end{quote}

학생은 이후 각 캐릭터를 고른 이유를 다음과 같이 설명했다. 아래는 학생의 설명 원문이다.
\par\noindent\begin{minipage}{\linewidth}
\textbf{카멜리아}
\begin{quote}
``일단 카멜리아 같은 경우, 가장 '밈'적으로 활용도가 높은 캐릭터이기 때문에 다양한 그림체가 있을 거라 판단했다. 그래서 그런 변형된 이미지도 잘 확인이 될 지 궁금했다.''
\end{quote}
\end{minipage}\par\medskip
\par\noindent\begin{minipage}{\linewidth}
\textbf{에이메스}
\begin{quote}
``에이메스는 내가 개인적으로 가장 좋아하는 캐릭터이기 때문에 골랐다. 별 다른 특별한 이유는 없다.''
\end{quote}
\end{minipage}\par\medskip
\par\noindent\begin{minipage}{\linewidth}
\textbf{양양}
\begin{quote}
``양양은 명조 게임의 마스코트라고 볼 수 있는 캐릭터이다. 게임 앱의 아이콘을 담당하기도 하기 때문에 게임을 하는 사람 입장에서는 의미가 깊은 캐릭터다. 또한 양양은 원래 4성 캐릭터이지만, 최근에 5성으로 업그레이드된 버전도 있기에 외형이 달라졌는데, 외형이 달라졌는데도 같은 캐릭터로 인식할지가 궁금했다.''
\end{quote}
\end{minipage}\par\medskip
\par\noindent\begin{minipage}{\linewidth}
\textbf{린네}
\begin{quote}
``린네 또한 내가 개인적으로 좋아하는 캐릭터이다.''
\end{quote}
\end{minipage}\par\medskip

네 캐릭터의 이름은 제공된 ImageNet 범주 목록에 없었다.
같은 게임 안의 캐릭터라도 머리, 의상, 그림체와 표현 방식이 달라지는 이미지를 구분하는 과제로 정했다.

\subsection*{What I expected the model to use}
학습 전에 학생이 적은 예상은 다음과 같다.
\begin{quote}
``머리카락 색과 모양, 입고 있는 옷의 스타일링''
\end{quote}

\subsection*{Images per category}
처음에는 검색어 ``명조 양양'', ``명조 카멜리아'', ``명조 린네'', ``명조 에이메스''로 각각 150장씩 수집했다.
이 600장으로 정리 전 실험을 모두 마친 후 정리했다.
린네는 최초 정리 후 43장이 남아 ``Lynae Wuthering Waves''로 50장을 추가 수집했고 그중 11장을 유지했다.
이 단계의 전체 다운로드 원본은 650장, 정리본은 311장이었으며 모델 tag는 \texttt{final}이었다.
이 수치는 아래에서 과거 실험의 결과로 보존한다.

학생이 ``너무 적어 좀 더 사진 찾아봐''라고 요청해 네 범주를 다시 수집했다.
영어·한국어·일본어·중국어의 캐릭터별 검색어를 사용했고, 제공된 수집 도구와 고정된 패키지 버전을 유지했다.
모든 새 후보를 시각 검토하고 기존 다운로드 원본, 다른 새 후보와 개인 평가용 사진에 대한 중복도 확인했다.
색이나 그림체가 비슷하다는 이유만으로 다른 캐릭터를 포함하지 않았다.
다운로드 사본은 원본 폴더에 모두 추가 보관하고, 제외 파일은 공식 제거 도구로 이동·기록했다.
현재 학습용 자료는 다음 표와 같다. 개인 촬영/캡처 24장은 이 자료에 포함하지 않았다.
\begin{center}\small
\begin{tabular}{lrrrrr}
\toprule
범주 & 최초 수집 & 1차 추가 & 이번 추가 & 전체 제거 & 현재 유지 \\
\midrule
린네 & 150 & 50 & 600 & 617 & 183 \\
양양 & 150 & 0 & 470 & 464 & 156 \\
에이메스 & 150 & 0 & 360 & 286 & 224 \\
카멜리아 & 150 & 0 & 240 & 210 & 180 \\
\bottomrule
\end{tabular}
\end{center}
현재 다운로드 원본은 2320장, 유지한 학습용 정리본은 743장이다.

이번 추가 수집·정리의 근거와 현재 클래스별 수는 아래에 보존한다.
\outputfile{outputs/expanded_summary.txt}

최초 수집과 린네 1차 추가 수집 출력도 역사 기록으로 그대로 보존한다.
\outputfile{outputs/collect.txt}
\outputfile{outputs/collect_supplement.txt}

\section*{Problem 2. Preparing an image}
\texttt{prepare\_image}는 RGB 변환, 짧은 변 256으로 bilinear 리사이즈,
중앙 224 정사각형 크롭, 픽셀 스케일링과 채널별 정규화, 채널 우선 배열 변환 순서로 구현했다.
반환값은 \texttt{float32}, 크기는 \texttt{(3, 224, 224)}다.
투명 PNG, 회색조, 대문자 확장자와 읽을 수 없는 파일을 처리했고,
각 범주가 학습과 테스트에 모두 들어가도록 분할했다.

\sidebyside{figures/prepare_original.jpg}{다운로드 원본: 에이메스 \texttt{0033.jpg}}
           {figures/prepare_result.png}{정규화를 되돌려 표시한 실제 모델 입력}
원본의 길게 뻗은 머리카락, 어깨 아래 의상과 왼쪽의 다른 인물은 중앙 크롭에서 잘린다.
얼굴과 머리 장식은 남지만, 원본 전체의 포즈와 옷의 형태를 모두 볼 수는 없다.
따라서 원본에서 보이는 특징과 모델 입력에서 보이는 특징을 구분해야 한다.

\section*{Problem 3. First training runs}
첫 네 실험은 정리 전에 동일한 600장과 동일한 분할로 실행했다.
모든 실험의 batch size는 8, seed는 0이다.
첫 네 실험은 480장으로 학습하고 120장으로 테스트했다.
Apple Silicon의 MPS를 사용했다. 추가 설정은 학습률을 바꾼 \texttt{lr3e4}이다.
기본 학습률 0.0001 대신 0.0003을 사용했고, 나머지 표의 조건은 같다.
이 최초 \texttt{lr3e4} 설정을 학생이 당시 직접 선택했는지는 대화 기록으로 확인되지 않았다.
아래의 \texttt{student\_lr3e4}는 이후 학생이 학습률 증가를 직접 고른 뒤 수행한 별도 정리 후 실험이다.
새 선택을 최초 정리 전 실험의 학생 선택으로 소급하여 기록하지 않는다.
\texttt{clean}은 최초 자료 정리 후, \texttt{final}은 린네 1차 추가 수집과 정리 후의 과거 모델이다.
이번 네 범주 추가 수집·정리 후의 모델은 \texttt{expanded}이며 현재 데모에서 사용하는 모델이다.
\begin{center}\scriptsize
\setlength{\tabcolsep}{3pt}
\begin{tabular}{lllrrrrr}
\toprule
tag & 시작 & 학습 층 & epochs & lr & 학습 파라미터 & 학습 정확도 & 테스트 정확도 \\
\midrule
\multicolumn{8}{l}{\textbf{정리 전 비교: 같은 600장, 학습 480장 / 테스트 120장}} \\
run & ImageNet & all & 10 & 0.0001 & 11,178,564 & 99.8\% & 78.3\% \\
scratch & random & all & 10 & 0.0001 & 11,178,564 & 99.4\% & 51.7\% \\
frozen & ImageNet & last layer & 10 & 0.001 & 2,052 & 89.2\% & 74.2\% \\
lr3e4 & ImageNet & all & 10 & 0.0003 & 11,178,564 & 99.8\% & 80.8\% \\
\midrule
\multicolumn{8}{l}{\textbf{자료 변경 후 기록: 정리 및 추가 수집으로 테스트 분할이 달라짐}} \\
clean & ImageNet & all & 10 & 0.0001 & 11,178,564 & 100.0\% & 88.5\% \\
final & ImageNet & all & 10 & 0.0001 & 11,178,564 & 100.0\% & 93.7\% \\
\midrule
\multicolumn{8}{l}{\textbf{현재 자료의 학습률 비교: 같은 743장, 동일 테스트 149장}} \\
expanded & ImageNet & all layers & 10 & 0.0001 & 11,178,564 & 100.0\% & 93.3\% \\
student\_lr3e4 & ImageNet & all layers & 10 & 0.0003 & 11,178,564 & 99.8\% & 93.3\% \\
\bottomrule
\end{tabular}
\end{center}
정리 전 첫 네 실행은 같은 테스트 파일에서 초기 가중치, 학습 층과 학습률에 따른 결과를 비교했다.
\texttt{clean}과 \texttt{final}은 자료와 분할이 달라 그 수치의 차이를 동일한 테스트 사진에서의 개선으로 해석하지 않는다.
\texttt{expanded}와 \texttt{student\_lr3e4}는 같은 자료와 동일한 순서의 테스트 파일에서 학습률을 비교한 정리 후 실험이다.
이 비교는 각 설정을 한 번씩 실행한 결과이며, 최초 정리 전 설정의 학생 선택 여부를 바꾸지 않는다.
\subsection*{Loss and generalization}
기본 실험의 학습 손실은 초반에 빠르게 내려가고 이후 낮은 수준에서 변동했다.
반면 테스트 손실은 초반 이후 꾸준히 개선되지 않았고 후반에 증가했다.
학습 정확도 99.8\%와 테스트 정확도 78.3\%의 차이는 21.5 percentage points다.
이는 학습 이미지에 잘 맞춘 결과가 같은 비율로 새로운 이미지까지 이어지지 않음을 보여준다.
정리 전에는 다른 캐릭터, 정보 화면과 근접 중복이 섞여 있어 숫자를 순수한 캐릭터 구분 능력으로 해석하기 어렵다.

기본 실험에서 가장 많이 혼동한 쌍은 린네와 양양이었다.
린네를 양양으로 5번, 양양을 린네로 3번 분류했다.
머리카락 색과 모양을 단서로 삼을 것이라는 최초 예상에 비추어 이 혼동을 확인했다.
배경, 다른 인물과 검색 라벨 오염도 함께 작용했을 가능성이 있다.

이 혼동 결과를 보고 학생은 다음과 같이 반응했다.
\begin{quote}
``린네와 양양을 혼동한 결과는 놀라웠다. 일단 사용하는 color palette 가 너무 다르기 때문이다. 양양하면 '파란색'이고, 린네하면 '네온 초록'과 '노랑' 계열인데, 어떻게 혼동될까 싶었다.''
\end{quote}

ImageNet에서 시작한 테스트 정확도는 78.3\%, 무작위 초기값에서 시작한 정확도는 51.7\%로 26.6 percentage points 차이였다.
작은 자료에서 처음부터 시각 특징을 배우는 것보다 이미 학습한 특징을 조정하는 편이 유리했던 결과로 해석한다.
마지막 층만 학습한 결과는 74.2\%로 전체 층 학습보다 4.1 percentage points 낮았다.
학습되는 파라미터가 2,052개여도 기본 분류에 상당 부분 도움이 되었으나, 전체 특징을 조정한 수준에는 도달하지 않았다.
학습률을 바꾼 실험은 80.8\%로 기본 실험보다 2.5 percentage points 높았다.

\subsection*{Figures for every run}
각 실험의 손실/정확도 곡선, 혼동 행렬과 자신 있게 틀린 테스트 이미지를 함께 보인다.
\texttt{clean}, \texttt{final}, \texttt{expanded}는 자료와 테스트 분할이 바뀌었으므로 최초 실험과 같은 테스트 집합의 성능 비교는 아니다.

\subsection*{run}
\par\noindent \picturefile{figures/run_curves.png}\par
\begin{center}\par\noindent \picturefile[0.62\linewidth]{figures/run_confusion.png}\par\end{center}
\par\noindent \picturefile{figures/run_worst.png}\par
\clearpage
\subsection*{scratch}
\par\noindent \picturefile{figures/scratch_curves.png}\par
\begin{center}\par\noindent \picturefile[0.62\linewidth]{figures/scratch_confusion.png}\par\end{center}
\par\noindent \picturefile{figures/scratch_worst.png}\par
\clearpage
\subsection*{frozen}
\par\noindent \picturefile{figures/frozen_curves.png}\par
\begin{center}\par\noindent \picturefile[0.62\linewidth]{figures/frozen_confusion.png}\par\end{center}
\par\noindent \picturefile{figures/frozen_worst.png}\par
\clearpage
\subsection*{lr3e4}
\par\noindent \picturefile{figures/lr3e4_curves.png}\par
\begin{center}\par\noindent \picturefile[0.62\linewidth]{figures/lr3e4_confusion.png}\par\end{center}
\par\noindent \picturefile{figures/lr3e4_worst.png}\par
\clearpage
\subsection*{clean}
\par\noindent \picturefile{figures/clean_curves.png}\par
\begin{center}\par\noindent \picturefile[0.62\linewidth]{figures/clean_confusion.png}\par\end{center}
\par\noindent \picturefile{figures/clean_worst.png}\par
\clearpage
\subsection*{final}
\par\noindent \picturefile{figures/final_curves.png}\par
\begin{center}\par\noindent \picturefile[0.62\linewidth]{figures/final_confusion.png}\par\end{center}
\par\noindent \picturefile{figures/final_worst.png}\par
\clearpage
\subsection*{expanded: current model}
현재 정리본 743장 중 594장은 학습에, 149장은 다운로드 테스트에 사용했다. ImageNet 초기값에서 all layers를 학습했다. epochs 10, lr 0.0001, batch size 8, seed 0이며, 학습 정확도는 100.0\%, 다운로드 테스트 정확도는 93.3\%다.
\par\noindent \picturefile{figures/expanded_curves.png}\par
\begin{center}\par\noindent \picturefile[0.62\linewidth]{figures/expanded_confusion.png}\par
\end{center}
\par\noindent \picturefile{figures/expanded_worst.png}\par
\clearpage
\subsection*{Student-selected learning-rate follow-up after cleaning}
학생은 추가 실험의 방향과 이유를 다음과 같이 직접 설명했다.
\begin{quote}
``학습 횟수보다는 학습률을 올리면 좋을 것 같아. 간혹 모델들은 학습을 반복해서 하면 좀 능률이 떨어지는 경우가 많은 것 같더라고''
\end{quote}

학습률을 올리는 방향은 학생이 선택했고, 실행 수치 0.0001에서 0.0003으로의 변경은 에이전트가 안내한 뒤 적용했다.
현재 정리본을 그대로 사용하고 학습 횟수는 10, batch size는 8, seed는 0으로 유지했다.
새 tag는 \texttt{student\_lr3e4}다. 기존 \texttt{expanded}와 같은 743장 중 594장으로 학습하고,
동일한 순서의 149장 테스트 파일로 평가했음을 저장 설정과 사전·사후 검증으로 확인했다.
ImageNet 가중치에서 새로 시작해 모든 층을 학습했으며, 기존 학습 결과를 이어서 학습한 것은 아니다.
분할과 모델 초기화에 같은 seed를 사용했다.

이는 학생의 이번 선택 이후 수행한 정리 후 추가 비교다.
기존 600장으로 수행했던 정리 전 \texttt{lr3e4}를 학생이 당시 선택했다고 바꾸지 않는다.
README의 본인 선택을 정리 전에 수행하라는 시간 순서를 과거에 충족했다고 주장하지 않는다.
실행 명령과 학생 원문, 두 실행의 설정 및 측정값은 \path{outputs/student_lr3e4_analysis.json}에 보존했다.

기존 실행의 학습 정확도는 100.0\%, 다운로드 테스트 정확도는 93.3\%였다.
새 실행은 각각 99.8\%와 93.3\%였고, 표시된 테스트 정확도의 차이는 0.0 percentage points다.
같은 개인 평가 사진에서 기존 실행은 22/24 (91.7\%),
새 실행은 22/24 (91.7\%)였다.

개인 평가 비교는 기존 \texttt{check.py}의 독립 전처리로 수행했다.
다운로드 테스트 수치는 각 실행 때 측정하여 저장한 값이며, 개인 사진 평가는 같은 24장에 대한 새 비교다.
동일 자료·설정의 한 번씩의 MPS 실행이므로 실행 변동과 작은 테스트 표본을 고려해야 한다.
학습률의 일반적인 우열을 확정하거나 반복 횟수에 따른 효율을 검증한 결과로 해석하지 않는다.
이번에는 반복 횟수를 바꾸지 않았으므로 학생의 반복 학습에 대한 생각은 선택 이유와 예상으로 기록한다.
새 모델은 추가 비교용으로 저장했다. 현재 웹 모델과 앞서 보고한 개인 평가 및 크롭 실험의 모델은 \texttt{expanded}다.

\outputfile{outputs/student_lr3e4_analysis.txt}
\outputfile{outputs/student_lr3e4_verification.txt}
\clearpage
\subsection*{Student-selected run: measured curves and mistakes}
\par\noindent\picturefile{figures/student_lr3e4_curves.png}\par
\begin{center}\picturefile[0.62\linewidth]{figures/student_lr3e4_confusion.png}\end{center}
\par\noindent\picturefile{figures/student_lr3e4_worst.png}\par
\clearpage
\subsection*{Student-selected run: complete printed training and checker outputs}
\outputfile{outputs/student_lr3e4_output.txt}
\outputfile{outputs/student_lr3e4_compare.txt}

\subsection*{Output of the runs}
\outputfile{outputs/run_output.txt}
\outputfile{outputs/scratch_output.txt}
\outputfile{outputs/frozen_output.txt}
\outputfile{outputs/lr3e4_output.txt}
\outputfile{outputs/clean_output.txt}
\outputfile{outputs/final_output.txt}
\outputfile{outputs/expanded_output.txt}
\section*{Problem 4. Cleaning}
\subsection*{Student rule stated after the existing runs}
학생은 기존 수집·정리·학습 결과를 확인한 뒤 다음 기준을 직접 제시했다.
\begin{quote}
``사진 정리 규칙은 일단 내가 말한 그 캐릭터의 특징 (머리 색, 입고 있는 옷) 등이 어떤 식으로든 clear 하게 나와야 한다. 또한 화질이 너무 떨어지는 건 안 된다''
\end{quote}
이는 이번에 받은 학생의 기준이며, 최초 정리 전에 학생이 적었던 규칙으로 소급하여 기록하지 않는다.
학생은 현재 정리본 검토에 대해 다음과 같이 확인했다.
\begin{quote}
``사진 정리 이미 다 확인했어. clean 에 있는 게 다 정리된 거니까.''
\end{quote}
현재 \path{data/clean/}의 사진 전체를 학생이 검토하고 유지하는 자료로 확인한 기록이다.
이미 제거한 후보의 과거 장별 판단까지 학생이 직접 수행했다고 소급하지 않는다.
화질의 구체적 수치 기준은 학생이 제시하지 않았으므로 임의로 정하지 않았다.
아래의 기존 제거 목록과 학습 결과는 에이전트가 적용했던 당시 기준의 기록이다.

\subsection*{Earlier agent rule and decisions}
대상 캐릭터를 식별할 수 있고 그 캐릭터가 주된 대상인 게임 화면과 공식/팬 일러스트를 유지했다.
다른 캐릭터, 글/무기/UI가 주된 이미지, 심하게 가려지거나 잘린 이미지는 제외하고, 같은 원본의 근접 중복은 올바른 클래스에서 명확한 대표본만 유지했다.

범주와 이미지 범위는 학생이 결정했다. 다음 진행 위임에 따라 세부 규칙과 장별 검토는 에이전트가 수행했다.
학생이 각 파일의 유지/제외를 직접 지정한 것처럼 기록하지 않는다.
\begin{quote}
``여기에도 정보가 있을거야 이것도 참고해주고, 지금부터는 short report 제외하고 전부 끝내놔. 커밋은 내가 할 거고 다 허용이니까 물어보지마''
\end{quote}
에이전트의 당시 규칙은 제거 전에 기록했고, 최초 600장을 모두 시트와 필요 시 원본 확대로 확인했다.
실제 유지 목록에는 머리/의상으로 대상 캐릭터를 식별할 수 있는 코스프레 사진도 일부 있다.
예를 들어 린네의 \texttt{0057.jpg}, \texttt{0124.jpg}, \texttt{0133.jpg}와 양양의 \texttt{0092.jpg}다.
이들은 최초 규칙에 적은 게임/일러스트라는 표현보다 범위를 넓혀 유지한 예외이며,
장별 판단 기록에 남겼다. 모든 정리본이 게임 화면과 그림만으로 구성됐다고 서술하지 않는다.
제거는 \texttt{clean.py remove}로 이동/기록했고 다운로드 원본을 수정하지 않았다.
기록의 공통 제거 사유가 잘못 분류된 세 항목은 장별 관찰 근거에 맞게 수정했다. 파일 위치와 학습 자료는 바뀌지 않았다.
\outputfile{outputs/cleaning_rule.txt}
\outputfile{outputs/clean_reason_correction.txt}

\subsection*{Three kinds of junk}
\par\noindent\begin{minipage}{\linewidth}
다른 캐릭터가 검색 라벨에 섞인 예: 카멜리아 폴더의 \texttt{0058.jpg}는 방랑자 화면이다.
\begin{center}\par\noindent \picturefile[0.55\linewidth]{figures/junk_other_character.jpg}\par\end{center}
\end{minipage}\par\medskip
\par\noindent\begin{minipage}{\linewidth}
무기 옵션과 글이 주된 정보 화면의 예: 에이메스 폴더의 \texttt{0030.jpg}.
\begin{center}\par\noindent \picturefile[0.4\linewidth]{figures/junk_weapon.jpg}\par\end{center}
\end{minipage}\par\medskip
\par\noindent\begin{minipage}{\linewidth}
작은 아이콘/글이 주된 표의 예: 린네 폴더의 \texttt{0055.jpg}.
\begin{center}\par\noindent \picturefile[0.55\linewidth]{figures/junk_table.jpg}\par\end{center}
\end{minipage}\par\medskip

\subsection*{What I removed}
아래는 최초 자료 정리와 린네 1차 추가 수집까지 반영한 과거 결과다. 이번 확장 결과는 별도로 구분한다.
최초 600장 중 300장을 유지했고, 린네 추가분 50장 중 11장을 유지했다.
전체 장별 판단은 ZIP의 \texttt{outputs/cleaning\_decisions.json},
\texttt{supplement\_decisions.json}와 이동 로그 \texttt{removed\_log.csv}에 있다.
\outputfile{outputs/clean_count_initial.txt}
\outputfile{outputs/clean_count.txt}
클래스별 제거 이유와 수는 아래 출력과 같다.
\begin{center}\small
\begin{tabular}{lrrrrrr}
\toprule 범주 & 근접 중복 & 다른 캐릭터 & 주 대상 아님 & 글/UI & 가림/크롭 & 합계 \\
\midrule
린네 & 67 & 61 & 7 & 6 & 5 & 146 \\
양양 & 18 & 17 & 53 & 1 & 2 & 91 \\
에이메스 & 39 & 9 & 6 & 6 & 0 & 60 \\
카멜리아 & 8 & 16 & 9 & 8 & 1 & 42 \\
\bottomrule
\end{tabular}\end{center}
\outputfile{outputs/clean_reason_by_class.txt}

도구가 처음 찾은 근접 중복 후보는 48쌍이다.
최종 판단 후 두 파일을 모두 유지한 후보쌍은 없었다.
도구가 찾지 않은 크롭/배너/반전 변형도 원본을 비교해 정리했다.
린네 1차 추가까지의 과거 제거 사유 중 근접 중복은 132개 파일이다. 후보쌍 수와 제거한 파일 수는 서로 다른 수치다.
아래 그림은 도구가 시트에 그린 후보이며 전체 쌍은 \texttt{outputs/suspects.csv}에 보존했다.
한글 글꼴이 없는 시트는 저장된 CSV의 이미지 순서와 점수를 그대로 읽어 글꼴만 바꿔 다시 그렸다.
확률과 유사도는 CSV에 저장된 소수 네 자리 문자열을 그대로 표시하며, 점수를 다시 측정하지 않았다.
\outputfile{outputs/cleaning_figure_render.txt}
\par\noindent \picturefile{figures/copies.png}\par

각 클래스 상위 20개 suspect 중 제거한 수는 다음과 같다.
\begin{center}\begin{tabular}{lr}
\toprule 범주 & 제거/검토 \\
\midrule

린네 & 19/20 \\
양양 & 18/20 \\
에이메스 & 6/20 \\
카멜리아 & 9/20 \\
\bottomrule\end{tabular}\end{center}
suspect 순위만으로 제거하지 않았다. 메카 표현이나 다른 그림체의 팬 일러스트처럼 학습 범위에 맞는 드문 표현은 유지했다.
상위 20개 밖에도 최초 자료에서 248장을 제거했다.
자동 순위는 전체 육안 검토를 대신할 수 없었다.
도구를 실행한 뒤 놓친 오염을 새로 발견했다고 가정하지 않고,
이미 검토한 오염과 순위가 얼마나 겹치는지 확인했다.

\par\noindent \picturefile{figures/suspects_lynae.png}\par
\par\noindent \picturefile{figures/suspects_yangyang.png}\par
\par\noindent \picturefile{figures/suspects_aemeath.png}\par
\par\noindent \picturefile{figures/suspects_camellia.png}\par
\subsection*{Before and after cleaning}
\outputfile{outputs/compare.txt}
정리 후 다운로드 테스트 정확도는 78.3\%에서 88.5\%로 높아졌지만,
동일한 개인 촬영 사진에서는 87.5\%에서 75.0\%로 낮아졌다.
다운로드 테스트 집합 자체가 바뀌었으므로 첫 숫자의 상승을 동일 이미지에서의 개선으로 보지 않는다.
새 출처 사진에서는 정리된 일러스트/게임 화면에 더 잘 맞춘 특징이 실제 행사장 배경, 입체 인형 탈과 다른 촬영 시점에 그대로 적용되지 않았을 가능성이 있다.
자료가 줄고 테스트 표본이 작아진 영향도 있다. 이 결과로 정리가 항상 새 출처 정확도를 높인다고 주장할 수 없다.

린네 1차 추가 수집까지 반영한 과거 \texttt{final}은 다운로드 테스트 93.7\%, 새 출처 79.2\%다.
정리와 추가 수집의 효과를 하나로 합쳐서 설명하지 않는다.
이 단계에서는 모든 범주의 정리본을 확보한 \texttt{final}을 모델로 지정했고,
새 출처 정확도가 높은 정리 전 모델로 바꿔 고르지 않았다.
모든 비교에서 개인 촬영 사진은 학습에 사용하지 않았다.

\subsection*{Additional collection and cleaning: expanded stage}
기존 정리 결과와 여섯 학습 모델은 보존하고, 이번에 수집한 후보에 동일한 장별 검토를 적용했다.
근접 중복 점수는 확인할 후보를 찾는 데 사용했다. 다른 포즈·표정·동작인 게임 프레임이나 SD 그림은 실제 이미지를 비교한 뒤 판단했다.
같은 원본의 단순 크롭·재압축·배경·워터마크 변경은 제외했다.
정확한 복장·머리 장식으로 식별되는 단독 코스프레는 앞서 명시한 예외 범위에서 유지했다.
개인 평가용 사진은 새 후보의 라벨 정리나 학습에 넣지 않았다.
전체 장별 유지/제외 사유와 검색 URL·파일 해시는 ZIP의 \path{outputs/expanded_decisions.json},
\path{outputs/expanded_import.json}와 \path{outputs/expanded_removed_log.csv}에 보존했다.
린네 이번 수집의 원래 터미널 출력은 파일로 저장되지 않았다.
ZIP의 \path{outputs/expanded_collect_lynae.txt}는 저장된 후보 파일과 SHA256을 다시 검사한 사후 검증 출력이며,
수집 당시 출력의 재구성이 아니다. 검색 출처 기록과 검색어별 실제 수집 요약은 별도 JSON에 보존했다.
아래는 실제 폴더에 추가하기 전에 유지 후보 전체를 기존 정리본·개인 평가 사진과 함께 검사한 기록이다.
이 출력의 ``repository data/results have not been changed''는 해당 사전 검사 시점을 가리킨다.
그 후 원본 추가와 공식 제거를 실행했으며, 현재 파일 수는 이어지는 정리 도구 출력에 보존한다.
\outputfile{outputs/expanded_collection_audit.txt}

공식 정리 도구의 현재 수는 아래와 같다.
\outputfile{outputs/expanded_clean_count.txt}

\texttt{final}과 \texttt{expanded}를 동일한 개인 촬영/캡처 24장으로 비교했다.
평가 사진의 정답률을 보고 과거 모델로 되돌려 고르지 않고, 추가 수집·정리를 완료한 \texttt{expanded}를 현재 모델로 정했다.
다운로드 테스트 분할이 바뀌므로 그 정확도의 차이를 같은 사진에서의 개선이라고 해석하지 않는다.
\outputfile{outputs/expanded_compare.txt}

\section*{Problem 5. Images from a new source}

\begingroup\raggedright
새 출처 평가용 이미지를 준비할 수 있는지에 대한 답변은 다음과 같다.
\begin{quote}
``모두 준비할 수 있고, 게임 이벤트에 가서 찍은 현수막 속의 사진 그리고 인형 탈 등도 있어''
\end{quote}

양양 이미지는 본인이 직접 촬영한 행사 사진과 직접 캡처한 인게임 화면이다.
행사 사진의 원본 파일은 \path{IMG_0804.HEIC}, \path{IMG_0809.HEIC},
\path{IMG_0810.HEIC}, \path{IMG_0812.HEIC}이며, 인게임 화면은
\path{IMG_1390.JPG}, \path{IMG_1392.JPG}이다.
캐릭터와 출처는 다음 답변으로 확인했다.
\begin{quote}
``모두 양양이며, 모두 직접 촬영·캡처했어요''
\end{quote}

린네 이미지로 제공한 행사 사진의 원본 파일은
\path{IMG_0844.HEIC}, \path{IMG_0845.HEIC}, \path{IMG_0850.jpg},
\path{IMG_0860.HEIC}, \path{IMG_0871.HEIC}이며, 인게임 화면은
\path{IMG_1417.JPG}이다. 출처는 다음 답변으로 확인했다.
\begin{quote}
``네, 모두 직접 촬영·캡처했어요''
\end{quote}
\path{IMG_0845.HEIC}에는 다른 인형 탈이 중앙에 크게 보이고 린네는 왼쪽에 있다.
학생이 지정한 라벨을 유지해 평가에 포함했다. 중앙 크롭에서는 다른 캐릭터가 주된 대상이 되므로 오분류 해석 시 이 모호성을 구분한다.

에이메스 이미지의 원본 파일은 \path{IMG_0553.HEIC}, \path{IMG_0802.HEIC},
\path{IMG_0806.HEIC}, \path{IMG_0818.HEIC}, \path{IMG_0869.HEIC},
\path{IMG_0215.PNG}이다. 캐릭터와 출처는 다음 말로 확인했다.
\begin{quote}
``이건 에이메스고 전부 내가 직접 찍은 것들이야.''
\end{quote}
\path{IMG_0553.HEIC}는 캐릭터 그림과 시계가 표시된 화면을 촬영한 사진이다.
\path{IMG_0802.HEIC}, \path{IMG_0806.HEIC}, \path{IMG_0818.HEIC}는
행사의 외부 현수막, 내부 현수막, 모니터를 각각 촬영한 사진이다.
\path{IMG_0869.HEIC}는 캐릭터 일러스트로 장식된 행사 차량을 촬영한 사진이며,
\path{IMG_0215.PNG}는 인게임 캡처 화면이다.
현수막과 모니터 사진에는 여러 캐릭터와 같은 일러스트가 보인다.
직접 촬영한 사진의 시점과 배경은 다르지만 일부는 같은 행사 일러스트를 담으므로 서로 완전히 독립인 표본으로 보기 어렵다. 평가에서 그대로 유지하고 이 상관관계를 한계로 기록한다.

카멜리아로 제공받은 파일은 \path{IMG_0867.HEIC}, \path{IMG_0834.HEIC},
\path{IMG_0583.jpg}, \path{IMG_0584.jpg}, \path{IMG_2554.jpg},
\path{IMG_2555.jpg}이다. 직접 촬영한 사진인지에 대한 답변은 다음과 같다.
\begin{quote}
``네, 모두 직접 촬영했어요''
\end{quote}
\path{IMG_0583.jpg}, \path{IMG_0584.jpg}, \path{IMG_0834.HEIC}는
행사장의 인형 탈을 측면 또는 정면에서 촬영한 사진이다.
\path{IMG_0867.HEIC}에는 인형 탈과 코스플레이 참가자가 함께 보인다.
\path{IMG_2554.jpg}는 진행자와 캐릭터 일러스트가 표시된 무대 화면을 촬영한 사진이며,
\path{IMG_2555.jpg}는 캐릭터 일러스트를 가까이 촬영한 사진이다.
다운로드 이미지와의 근접 중복 여부는 overlap 도구로 검사했다.

HEIC 원본은 별도로 보관하고, 화면에 보이는 방향을 유지한 JPEG 사본으로 변환했다.
이 이미지들은 새 출처 평가용으로 보관하며 학습에 사용하지 않는다.


\endgroup
\subsection*{Sources and overlap}
각 범주에 6장씩, 전체 24장을 사용했다.
학생이 직접 촬영/캡처한 개인 파일이므로 다운로드한 인터넷 파일과 출처가 다르다.
실제 행사장의 조명, 인형 탈, 자동차와 주변 인물이 나타나 다운로드 일러스트/게임 화면과 표현 방식도 다르다.
파일 단위의 중복 검사와 특징 기반 overlap 검사 결과를 함께 확인했다.
아래의 과거 \texttt{final} 단계 overlap은 650개 원본 중 서로 다른 다운로드 이미지 649개와 비교한 결과다.
그 단계에서는 새 사진의 근접 중복이 없었고 가장 가까운 유사도는 0.86으로 기준 0.95보다 낮았다.
같은 행사 그림을 다른 위치에서 찍은 사진은 상관관계가 있으므로 24장을 완전히 독립적인 표본으로 가정하지 않는다.
\outputfile{outputs/new_images_received.txt}
\outputfile{outputs/overlap.txt}
확장 수집 후 전체 다운로드 원본과 같은 개인 사진 24장을 다시 검사한 현재 결과는 다음과 같다.
\outputfile{outputs/expanded_overlap.txt}

\subsection*{Downloaded test versus new images}
\begin{center}\begin{tabular}{lrr}
\toprule 모델 & 다운로드 테스트 & 직접 촬영/캡처 \\
\midrule

run & 78.3\% & 87.5\% \\
clean & 88.5\% & 75.0\% \\
final & 93.7\% & 79.2\% \\
expanded & 93.3\% & 91.7\% \\
\bottomrule\end{tabular}\end{center}
과거 \texttt{final}의 새 출처 결과는 19/24이었다. 현재 \texttt{expanded}의 결과는 22/24 (91.7\%)다.
인게임 캡처의 화면 표현은 익숙하지만, 행사 사진에는 무대 조명, 실물 원단의 질감,
카메라 시점, 군중, 차량, 휴대전화와 다른 캐릭터가 함께 나타난다.
캐릭터가 프레임에서 작거나 중앙에서 벗어나면 중앙 크롭이 핵심 특징을 제거하기도 한다.

\subsection*{Current expanded model: confident mistakes}
현재 모델은 개인 평가 사진에서 두 장만 틀렸다. 따라서 현재 모델의 오답을 다섯 장으로 채우지 않고 실제 두 사례를 모두 기록한다. 과거 모델의 다섯 오답은 같은 Problem 5에서 별도로 제시한다.
\par\noindent \picturefile{figures/expanded_new_worst.png}\par
\par\noindent\begin{minipage}{\linewidth}
\subsection*{1. IMG\_2554.jpg}
정답: 명조\_카멜리아; 예측: 명조\_에이메스 80.8\%
\begin{center}\par\noindent \picturefile[0.35\linewidth]{figures/expanded_new_mistake_1.jpg}\par
\end{center}
무대 오른쪽에 SD 캐릭터 그림이 있고 큰 행사 글씨와 진행자가 함께 보인다. 실제 중앙 입력에서는 캐릭터의 얼굴·의상 대부분이 오른쪽에서 잘리고 글씨와 진행자 일부가 크게 남는다. 전체 사진의 사용자 라벨과 중앙 입력의 주된 정보가 다를 수 있다.
\end{minipage}\par\medskip
\par\noindent\begin{minipage}{\linewidth}
\subsection*{2. IMG\_0860.jpg}
정답: 명조\_린네; 예측: 명조\_에이메스 45.6\%
\begin{center}\par\noindent \picturefile[0.45\linewidth]{figures/expanded_new_mistake_2.jpg}\par
\end{center}
입체 인형 탈의 큰 눈과 풍선껌이 중심에 보이고, 앞쪽 휴대전화·손과 옆의 다른 인형 탈 일부가 함께 나타난다. 실제 중앙 입력은 머리·얼굴 위주이며 아래쪽 의상은 크롭과 휴대전화로 대부분 가려진다. 게임 화면과 비율·질감·조명이 다르다.
위 설명은 관찰에 근거한 가능한 해석이며 특정 원인을 증명한 결과가 아니다. 점수는 네 범주 안의 상대 점수이며, 그 이미지가 범주에 적합하다는 확률을 보장하지 않는다.
\end{minipage}\par\medskip

\clearpage
\subsection*{Historical final model: five confident mistakes}
\par\noindent 아래 다섯 사례는 확장 전 \texttt{final} 모델에서 실제로 틀린 사진이다. 현재 모델의 오답과 구분하여 당시의 답과 점수를 그대로 보존한다.\par
\noindent \picturefile{figures/new_worst.png}\par

\par\noindent\begin{minipage}{\linewidth}
\subsection*{1. IMG\_2554.jpg}
정답: 명조\_카멜리아; 예측: 명조\_에이메스 89.6\%
\begin{center}\par\noindent \picturefile[0.45\linewidth]{figures/new_mistake_1.jpg}\par\end{center}
캐릭터가 무대 오른쪽에 작게 있고 가운데에는 진행자와 큰 글씨가 있다. 파란 화면과 조명, 중앙 크롭의 비캐릭터 정보가 결과에 영향을 주었을 가능성이 있다.
\end{minipage}\par\medskip

\par\noindent\begin{minipage}{\linewidth}
\subsection*{2. IMG\_0845.jpg}
정답: 명조\_린네; 예측: 명조\_양양 77.5\%
\begin{center}\par\noindent \picturefile[0.45\linewidth]{figures/new_mistake_2.jpg}\par\end{center}
린네는 왼쪽을 향한 뒤/측면이고 다른 인형 탈이 중앙에 크게 보인다. 학생이 지정한 전체 사진의 라벨과 모델의 중앙 입력이 보여 주는 주된 대상이 서로 다르다.
\end{minipage}\par\medskip

\par\noindent\begin{minipage}{\linewidth}
\subsection*{3. IMG\_0844.jpg}
정답: 명조\_린네; 예측: 명조\_에이메스 68.3\%
\begin{center}\par\noindent \picturefile[0.45\linewidth]{figures/new_mistake_3.jpg}\par\end{center}
린네 인형 탈의 옆모습이며 큰 휴대전화와 손이 아래를 가린다. 청록 풍선껌, 원단과 파란 무대 배경이 게임 모델과 다른 형태로 보인다.
\end{minipage}\par\medskip

\par\noindent\begin{minipage}{\linewidth}
\subsection*{4. IMG\_0860.jpg}
정답: 명조\_린네; 예측: 명조\_에이메스 58.8\%
\begin{center}\par\noindent \picturefile[0.45\linewidth]{figures/new_mistake_4.jpg}\par\end{center}
입체 인형 탈, 큰 눈과 풍선껌이 원래 게임 얼굴의 비율과 다르고, 손/휴대전화가 의상을 가린다. 중심 입력에서 익숙한 의상 정보가 약해졌을 가능성이 있다.
\end{minipage}\par\medskip

\par\noindent\begin{minipage}{\linewidth}
\subsection*{5. IMG\_0871.jpg}
정답: 명조\_린네; 예측: 명조\_에이메스 55.8\%
\begin{center}\par\noindent \picturefile[0.45\linewidth]{figures/new_mistake_5.jpg}\par\end{center}
차량 앞부분의 일러스트를 행사장에서 촬영한 사진이다. 중앙 크롭에 차량과 반사, 주변 배경이 크게 들어가고 원화의 얼굴/의상 정보는 작고 비스듬하다.
\end{minipage}\par\medskip
이 설명은 관찰에 근거한 가능한 해석이며 특정 원인을 이미 증명했다는 뜻은 아니다.
신뢰도는 네 범주 중 상대 점수이고, 이미지가 범주에 적합하다는 확률을 보장하지 않는다.



\section*{Problem 6. The demo}
현재 데모 모델은 추가 수집·정리 후 학습한 \texttt{expanded}이다.
같은 RGB 샘플, Pillow-compatible bilinear 필터와 half-even 중앙 크롭을 유지하고 이 체크포인트를 다시 내보냈다.
현재 모델의 Python/ONNX 비교 결과는 다음 출력과 같다.
\outputfile{outputs/expanded_export.txt}

기존 \texttt{final} 검증 때 발견한 단계적 축소와 ICC 해석 차이를 수정한 과정은 아래 과거 출력에 보존했다.
업로드 사본의 ICC 메타데이터는 메모리에서만 제거하며, 원본 디스크 사진은 바꾸지 않는다.
이번에는 \texttt{expanded}의 원본 개인 사진과 기존 변환본을 실제 브라우저에 다시 업로드해 확인했다.
\outputfile{outputs/browser_preprocessing.txt}

계획한 공개 주소는 \url{https://jellykeeng.github.io/mas1004-assignment1/},
코드 주소는 \url{https://github.com/jellykeeng/mas1004-assignment1}이다.
이 문서 작성 시점에는 공개 저장소 생성/게시 완료를 확인하지 못했다.
GitHub 인증 확인 명령의 실행 승인이 거부되어 계정 설정을 중단했고,
학생이 커밋을 직접 하겠다고 요청했으므로 에이전트는 커밋하거나 푸시하지 않았다.
따라서 위 주소는 배포 예정 주소이며, 제출 전 실제 공개 URL과 동작을 확인해야 한다.

범주에 없는 무기 사진을 현재 모델에 넣어도 ``모름'' 대신 명조\_린네 59.3\%를 출력했다.
이 모델에는 unknown 범주가 없어 다른 물체의 입력을 거부하지 못한다. 현재 초록색 self-test, 업로드 검사와 범주 밖 입력의 브라우저 화면은 아래에 보존한다.
\outputfile{outputs/expanded_browser_demo.txt}
\par\noindent \picturefile{figures/expanded_browser_selftest.png}\par
\par\noindent \picturefile{figures/expanded_browser_ood.png}\par

\section*{Problem 7. What the model looks at}
학생이 직접 선택한 전체 개인 사진의 머리카락 가림과 결과 해석은 \hyperref[sec:student-hair]{후속 머리카락 가림 절}에 기록했다.
\subsection*{Student-selected crops: results at a glance}
다음은 학생이 오분류 사진을 보고 직접 지정한 두 크롭과, 이후 직접 제공한 네 크롭본의 기존 결과를 모은 것이다.
모두 같은 \texttt{expanded} 모델과 기존 전처리를 사용했다.
\begin{center}\small
\setlength{\tabcolsep}{4pt}
\begin{tabular}{llrrl}
\toprule
범주 & 원본 파일 & 원본 정답 점수 & 변경 정답 점수 & 모델의 선택 \\
\midrule
카멜리아 & \texttt{IMG\_2554.jpg} & 0.9\% & 100.0\% & 교정: 에이메스$\to$카멜리아 \\
린네 & \texttt{IMG\_0860.jpg} & 17.2\% & 67.5\% & 교정: 에이메스$\to$린네 \\
린네 & \texttt{IMG\_0845.jpg} & 65.7\% & 87.5\% & 유지: 린네 \\
에이메스 & \texttt{IMG\_0806.jpg} & 69.6\% & 100.0\% & 유지: 에이메스 \\
에이메스 & \texttt{IMG\_0818.jpg} & 75.8\% & 100.0\% & 유지: 에이메스 \\
린네 & \texttt{IMG\_0844.jpg} & 80.9\% & 68.5\% & 유지: 린네 \\
\bottomrule
\end{tabular}
\end{center}
표의 점수는 정답 라벨의 softmax 점수이며, 가장 높은 점수로 선택된 답과 구분한다.
두 오답은 원본 0/2에서 크롭 후 2/2로 바뀌었다.
네 제공본은 원본과 변경본 모두 4/4를 맞혔으며 정답 점수는 세 사례에서 오르고 한 사례에서 낮아졌다.
이는 선택된 사진들의 전후 비교이며, 원래 개인 평가 22/24 (91.7\%)를 바꾼 새 정확도가 아니다.

학생은 카멜리아가 작게 보이고 린네의 옷 부분이 잘렸을 가능성을 근거로 캐릭터 중심/아래쪽 크롭을 요청했다.
그 관찰과 가설, 정확한 요청 원문은 \hyperref[p7:student_crop_request]{학생이 선택한 두 오답의 크롭} 절에 보존했다.
이후 네 제공본의 정확한 파일 대응과 전후 결과는 \hyperref[p7:student_provided_crops]{학생 제공 크롭 파일} 절에 있다.
\texttt{IMG\_0844}에서 정답 이름이 유지되면서 점수가 낮아진 결과는 학생의 예상과 달랐다.
학생은 결과를 확인한 뒤 중앙 정사각형 입력에서 의상이 일부 잘릴 수 있다는 잠정 추측을 제시했다.
이 사후 해석 원문과 실제 입력의 관찰은 \hyperref[p7:student_crop_reflection]{학생의 크롭 결과 해석} 절에서 구분해 기록한다.
배경, 캐릭터의 크기/위치와 보이는 얼굴/의상이 함께 바뀌어 하나의 특징만의 인과 효과를 확정하지 않는다.
아래의 과거 흑백/가림 실험은 에이전트 선택 기록이며, 학생의 직접 선택으로 바꾸어 서술하지 않는다.

학생의 학습 전 예상은 ``머리카락 색과 모양, 입고 있는 옷의 스타일링''이었다.
이 예상과 관련된 변화를 살펴보기 위해 흑백 변환, 머리/얼굴 영역 가림, 몸통/의상 영역 가림을 사용했다.
아래의 기존 흑백/가림 실험을 준비할 당시에는 학생이 구체적인 변환과 파일명을 직접 지정한 요청이 없었다.
이 세 실험의 구체적 변환 선택과 좌표는 진행 위임을 받은 에이전트가 결정한 것이며,
학생이 스스로 선택한 실험이라고 대신 서술하지 않는다.

관련된 학생의 요청을 그대로 적으면 다음과 같다.
\begin{quote}
``여기에도 정보가 있을거야 이것도 참고해주고, 지금부터는 short report 제외하고 전부 끝내놔. 커밋은 내가 할 거고 다 허용이니까 물어보지마''
\end{quote}

양양의 인게임 화면 두 장, 카멜리아의 인형 탈/SD 그림 사진,
린네의 인게임/인형 탈 사진, 에이메스의 인게임/화면 사진을 사용했다.
원본 파일은 그대로 두고, 모델 입력을 정확히 보존하는 중앙 256 정사각형 사본에만 변환을 적용해
\texttt{data/changed/}에 저장했다.
각 사본을 다시 준비한 배열과 원본을 준비한 배열이 픽셀 단위로 같음을 확인했다.
모델은 그 사본의 중앙 224 영역을 보므로 가림 위치가 실제 입력 안에 들어온다.
아래 ``원본'' 그림은 이 입력용 중앙 사본이다.
그림과 기존 23개 변환은 그대로 보존하고, 아래 캡션의 답과 점수만 현재 \texttt{expanded}로 다시 측정했다.
과거 \texttt{final} 예측은 ZIP의 \path{outputs/analysis.json}, 전체 분석 출력과 확장 전 보고서 원문에 보존했다.

\clearpage
\subsection*{Guess 1. Hair colour}
캐릭터별 머리색이 구분의 중요한 단서라는 최초 예상을 확인한다. 전체 중앙 사본을 흑백으로 바꿔 색 정보의 영향을 줄였다. 의상과 배경 색도 함께 바뀌므로 머리색만의 영향으로 단정할 수는 없다.
\par\noindent\begin{minipage}{\linewidth}
\texttt{IMG\_1390.JPG}
\sidebyside{figures/gray_01_original.png}{원본: 명조\_양양 99.5\%}{figures/gray_01_changed.png}{변환: 명조\_양양 99.9\%}
\end{minipage}\par\medskip
\par\noindent\begin{minipage}{\linewidth}
\texttt{IMG\_1392.JPG}
\sidebyside{figures/gray_02_original.png}{원본: 명조\_양양 99.8\%}{figures/gray_02_changed.png}{변환: 명조\_양양 99.2\%}
\end{minipage}\par\medskip
\par\noindent\begin{minipage}{\linewidth}
\texttt{IMG\_0584.jpg}
\sidebyside{figures/gray_03_original.png}{원본: 명조\_카멜리아 100.0\%}{figures/gray_03_changed.png}{변환: 명조\_에이메스 67.9\%}
\end{minipage}\par\medskip
\par\noindent\begin{minipage}{\linewidth}
\texttt{IMG\_2555.jpg}
\sidebyside{figures/gray_04_original.png}{원본: 명조\_카멜리아 100.0\%}{figures/gray_04_changed.png}{변환: 명조\_양양 63.9\%}
\end{minipage}\par\medskip
\par\noindent\begin{minipage}{\linewidth}
\texttt{IMG\_1417.JPG}
\sidebyside{figures/gray_05_original.png}{원본: 명조\_린네 96.6\%}{figures/gray_05_changed.png}{변환: 명조\_린네 90.7\%}
\end{minipage}\par\medskip
\par\noindent\begin{minipage}{\linewidth}
\texttt{IMG\_0850.jpg}
\sidebyside{figures/gray_06_original.png}{원본: 명조\_린네 87.5\%}{figures/gray_06_changed.png}{변환: 명조\_에이메스 82.5\%}
\end{minipage}\par\medskip
\par\noindent\begin{minipage}{\linewidth}
\texttt{IMG\_0215.PNG}
\sidebyside{figures/gray_07_original.png}{원본: 명조\_에이메스 100.0\%}{figures/gray_07_changed.png}{변환: 명조\_에이메스 76.9\%}
\end{minipage}\par\medskip
\par\noindent\begin{minipage}{\linewidth}
\texttt{IMG\_0553.jpg}
\sidebyside{figures/gray_08_original.png}{원본: 명조\_에이메스 99.9\%}{figures/gray_08_changed.png}{변환: 명조\_에이메스 93.2\%}
\end{minipage}\par\medskip
현재 \texttt{expanded}의 원본 8/8에서 변환 후 5/8를 맞췄고 예측 클래스가 3번 바뀌었다. 이는 선택한 사진에서 전체 색 정보를 제거한 뒤 관찰한 결과다. 머리색의 효과와 의상·배경·피부 색의 효과를 분리한 결과는 아니다.
\clearpage
\subsection*{Guess 2. Hair and head shape}
머리카락 모양과 머리 주변의 특징이 필요한지 확인한다. 표시한 머리/얼굴 영역을 RGB 128의 회색 사각형으로 가렸다. 머리카락뿐 아니라 얼굴과 장식도 함께 가려지는 조작이므로 머리 모양만 따로 분리한 실험은 아니다.
\par\noindent\begin{minipage}{\linewidth}
\texttt{IMG\_1390.JPG} 가림 좌표 [107, 20, 208, 110]
\sidebyside{figures/head_mask_01_original.png}{원본: 명조\_양양 99.5\%}{figures/head_mask_01_changed.png}{변환: 명조\_양양 92.3\%}
\end{minipage}\par\medskip
\par\noindent\begin{minipage}{\linewidth}
\texttt{IMG\_1392.JPG} 가림 좌표 [42, 22, 130, 93]
\sidebyside{figures/head_mask_02_original.png}{원본: 명조\_양양 99.8\%}{figures/head_mask_02_changed.png}{변환: 명조\_양양 97.9\%}
\end{minipage}\par\medskip
\par\noindent\begin{minipage}{\linewidth}
\texttt{IMG\_0584.jpg} 가림 좌표 [45, 16, 196, 118]
\sidebyside{figures/head_mask_03_original.png}{원본: 명조\_카멜리아 100.0\%}{figures/head_mask_03_changed.png}{변환: 명조\_카멜리아 100.0\%}
\end{minipage}\par\medskip
\par\noindent\begin{minipage}{\linewidth}
\texttt{IMG\_2555.jpg} 가림 좌표 [16, 16, 230, 140]
\sidebyside{figures/head_mask_04_original.png}{원본: 명조\_카멜리아 100.0\%}{figures/head_mask_04_changed.png}{변환: 명조\_카멜리아 100.0\%}
\end{minipage}\par\medskip
\par\noindent\begin{minipage}{\linewidth}
\texttt{IMG\_1417.JPG} 가림 좌표 [95, 25, 171, 89]
\sidebyside{figures/head_mask_05_original.png}{원본: 명조\_린네 96.6\%}{figures/head_mask_05_changed.png}{변환: 명조\_린네 83.0\%}
\end{minipage}\par\medskip
\par\noindent\begin{minipage}{\linewidth}
\texttt{IMG\_0850.jpg} 가림 좌표 [80, 56, 193, 150]
\sidebyside{figures/head_mask_06_original.png}{원본: 명조\_린네 87.5\%}{figures/head_mask_06_changed.png}{변환: 명조\_린네 50.9\%}
\end{minipage}\par\medskip
\par\noindent\begin{minipage}{\linewidth}
\texttt{IMG\_0215.PNG} 가림 좌표 [98, 16, 212, 115]
\sidebyside{figures/head_mask_07_original.png}{원본: 명조\_에이메스 100.0\%}{figures/head_mask_07_changed.png}{변환: 명조\_에이메스 99.8\%}
\end{minipage}\par\medskip
\par\noindent\begin{minipage}{\linewidth}
\texttt{IMG\_0553.jpg} 가림 좌표 [16, 16, 239, 192]
\sidebyside{figures/head_mask_08_original.png}{원본: 명조\_에이메스 99.9\%}{figures/head_mask_08_changed.png}{변환: 명조\_에이메스 39.8\%}
\end{minipage}\par\medskip
현재 \texttt{expanded}의 원본 8/8에서 변환 후 8/8를 맞췄고 예측 클래스가 0번 바뀌었다. 양양과 린네의 긴 머리는 가림 밖에 일부 남고, 에이메스의 얼굴 사진에서는 대부분의 시각 정보가 가려진다. 머리카락만의 효과를 분리할 수 없으며, 답이 유지됐다는 이유로 머리 특징을 사용하지 않는다고 단정할 수도 없다.
\clearpage
\subsection*{Guess 3. Outfit styling}
몸통과 의상 부분이 필요한지 확인한다. 그 영역을 회색 사각형으로 가렸다. 얼굴만 있는 에이메스 화면 사진은 가릴 옷이 없어 이 조작에서 제외했다. 사각형에는 일부 팔/피부도 포함되므로 의상만의 효과로 해석하지 않는다.
\par\noindent\begin{minipage}{\linewidth}
\texttt{IMG\_1390.JPG} 가림 좌표 [89, 135, 197, 239]
\sidebyside{figures/clothes_mask_01_original.png}{원본: 명조\_양양 99.5\%}{figures/clothes_mask_01_changed.png}{변환: 명조\_양양 99.5\%}
\end{minipage}\par\medskip
\par\noindent\begin{minipage}{\linewidth}
\texttt{IMG\_1392.JPG} 가림 좌표 [62, 133, 152, 239]
\sidebyside{figures/clothes_mask_02_original.png}{원본: 명조\_양양 99.8\%}{figures/clothes_mask_02_changed.png}{변환: 명조\_양양 99.3\%}
\end{minipage}\par\medskip
\par\noindent\begin{minipage}{\linewidth}
\texttt{IMG\_0584.jpg} 가림 좌표 [67, 127, 182, 238]
\sidebyside{figures/clothes_mask_03_original.png}{원본: 명조\_카멜리아 100.0\%}{figures/clothes_mask_03_changed.png}{변환: 명조\_카멜리아 98.6\%}
\end{minipage}\par\medskip
\par\noindent\begin{minipage}{\linewidth}
\texttt{IMG\_2555.jpg} 가림 좌표 [17, 145, 239, 239]
\sidebyside{figures/clothes_mask_04_original.png}{원본: 명조\_카멜리아 100.0\%}{figures/clothes_mask_04_changed.png}{변환: 명조\_카멜리아 97.3\%}
\end{minipage}\par\medskip
\par\noindent\begin{minipage}{\linewidth}
\texttt{IMG\_1417.JPG} 가림 좌표 [85, 120, 174, 230]
\sidebyside{figures/clothes_mask_05_original.png}{원본: 명조\_린네 96.6\%}{figures/clothes_mask_05_changed.png}{변환: 명조\_린네 94.5\%}
\end{minipage}\par\medskip
\par\noindent\begin{minipage}{\linewidth}
\texttt{IMG\_0850.jpg} 가림 좌표 [90, 148, 189, 239]
\sidebyside{figures/clothes_mask_06_original.png}{원본: 명조\_린네 87.5\%}{figures/clothes_mask_06_changed.png}{변환: 명조\_린네 96.9\%}
\end{minipage}\par\medskip
\par\noindent\begin{minipage}{\linewidth}
\texttt{IMG\_0215.PNG} 가림 좌표 [56, 133, 213, 239]
\sidebyside{figures/clothes_mask_07_original.png}{원본: 명조\_에이메스 100.0\%}{figures/clothes_mask_07_changed.png}{변환: 명조\_에이메스 100.0\%}
\end{minipage}\par\medskip
현재 \texttt{expanded}의 원본 7/7에서 변환 후 7/7를 맞췄고 예측 클래스가 0번 바뀌었다. 린네 인게임 사진의 사각형에는 맨살인 복부가, 인형 탈 사진의 사각형에는 아래쪽 얼굴/턱과 앞쪽 손도 포함된다. 의상 스타일만의 효과나 한 캐릭터 전체의 특성으로 일반화할 수 없다.
\subsection*{What the evidence supports}
현재 모델에 대해 같은 사진·변환을 다시 측정한 결과를 관찰했다.
학생의 최초 예상과 관련된 단서를 시험했지만, 머리색/모양과 의상 스타일에 대한 인과적 의존을 확정하지는 못했다.
흑백은 배경과 얼굴 색도 바꾸고, 사각형은 여러 특징을 동시에 제거하며 낯선 경계도 만든다.
원본들은 과거 \texttt{final}에서 정답이었던 작은 표본에서 선택했고, 현재 모델에 대해 새로 독립적으로 선정한 자료가 아니다.
가림 면적과 실험 분모도 서로 다르므로 변형 간 특징 중요도 순위를 매기거나 모든 사진으로 일반화하지 않는다.
세 변환은 에이전트가 최초 예상을 나누어 살펴본 과거 병렬 실험이다. 이 세 실험의 선택을 학생이 직접 수행한 반복 탐구로 소급하지 않는다. 이후 학생이 두 오답을 보고 직접 정한 크롭 가설과 요청은 아래에서 별도로 기록한다.
브라우저에도 같은 원본과 변환본을 업로드한 현재 검증을 아래에 보존한다.
\outputfile{outputs/expanded_browser_ablations.txt}


\clearpage
\subsection*{Student-selected follow-up: reframing two mistakes}
\label{p7:student_crop_request}
현재 모델이 직접 촬영한 사진 두 장을 에이메스로 오분류한 결과를 본 뒤,
학생은 각 사진의 구도에 대한 가설과 구체적인 크롭 방향을 직접 제시했다.
학생의 관찰과 가설을 원문으로 기록하면 다음과 같다.
\begin{quote}
``내 생각에 카멜리아 같은 경우는 사진이 너무 큰데 카멜리악 보이는 부분은 작아서 그런것 같아. 린네를 넣었는데 에이메스로 잘못 판단한 건 배경이 핑크색이기도 하고, 풍선껌 때문에 얼굴이 제대로 안 보이고, 옷 부분이 잘려서 린네인 걸 못 알아본 것 같아.''
\end{quote}
실행한 크롭은 다음 요청에 따른 것이다.
\begin{quote}
``카멜리아 같은 경우는 배경은 자르고, 카멜리아 부분만 사진을 크롭하면 맞출 수 있을 것 같고, 린네는 좀 첨부한 것처럼 좀 더 아래쪽에서 crop 하면 될 거 같아.''
\end{quote}
직전 대화에서 다룬 파일은 카멜리아의 \texttt{IMG\_2554.jpg}와
린네의 \texttt{IMG\_0860.jpg}다. 학생이 선택한 대상과 크롭 방향을
에이전트가 원본 픽셀 좌표로 옮겼으며, 좌표는 예측값을 확인하기 전에 고정했다.
린네 좌표는 첨부 예시의 구도를 참고한 것으로, 첨부 파일과 바이트가 동일하다고 주장하지 않는다.
기존 \texttt{expanded} 모델을 그대로 사용했고 새 학습이나 좌표별 예측 탐색은 하지 않았다.
원본과 평가 폴더를 그대로 보존하며, 정확한 RGB 픽셀 크롭의 PNG 사본만
\path{data/changed/student_crops/}에 저장했다. 기존 \texttt{prepare\_image}의
resize/중앙 크롭 절차도 바꾸지 않았다.

\begin{center}\begin{tabular}{lrr}
\toprule 사진 & 원본의 정답 점수 & 크롭 후 정답 점수 \\
\midrule
카멜리아 & 0.9\% & 100.0\% \\
린네 & 17.2\% & 67.5\% \\
\bottomrule\end{tabular}\end{center}
이 표의 점수는 각각 정답 라벨의 softmax 점수다.
모델이 선택한 이름과 그 이름의 점수는 아래 각 그림에 표시한다.
이 두 선택된 오답에서 원본은 0/2, 크롭 후에는 2/2를 맞혔다.
이는 선택된 두 사진의 전후 비교이며, 독립된 새 평가 자료의 정확도가 아니다.
원래 개인 평가 사진의 결과인 22/24 (91.7\%)는 그대로 보존한다.

\clearpage
\subsection*{Camellya: selected crop}
\texttt{IMG\_2554.jpg}의 크롭 좌표: \texttt{[820, 390, 1120, 702]}. 원본 크기 1290$\times$702에서 300$\times$312로 바뀌었다. 좌상단을 포함하고 우하단을 제외하는 원본 RGB 픽셀 좌표다.
\par\noindent\begin{minipage}{\linewidth}
\sidebyside{figures/expanded_new_mistake_1.jpg}{원본: 명조\_에이메스 80.8\%}{figures/student_crop_camellya_changed.png}{크롭: 명조\_카멜리아 100.0\%}
위는 전체 원본과 요청한 크롭 사본이다. 아래는 각각 모델에 실제로 들어간 224 정사각형 입력이다.
\sidebyside{figures/student_crop_camellya_before_input.png}{원본의 실제 입력}{figures/student_crop_camellya_after_input.png}{크롭 사본의 실제 입력}
카멜리아 그림을 중심으로 잘라 큰 행사 글씨와 진행자를 대부분 제외했다. 실제 중앙 입력에도 얼굴과 의상의 주된 부분이 들어온다. 대상의 위치와 크기, 배경과 주변 인물이 동시에 달라지므로 배경색만의 효과를 분리한 실험은 아니다.
\end{minipage}\par\medskip

\clearpage
\subsection*{Lynae: selected crop}
\texttt{IMG\_0860.jpg}의 크롭 좌표: \texttt{[0, 1248, 3024, 4032]}. 원본 크기 3024$\times$4032에서 3024$\times$2784로 바뀌었다. 좌상단을 포함하고 우하단을 제외하는 원본 RGB 픽셀 좌표다.
\par\noindent\begin{minipage}{\linewidth}
\sidebyside{figures/expanded_new_mistake_2.jpg}{원본: 명조\_에이메스 45.6\%}{figures/student_crop_lynae_changed.png}{크롭: 명조\_린네 67.5\%}
위는 전체 원본과 요청한 크롭 사본이다. 아래는 각각 모델에 실제로 들어간 224 정사각형 입력이다.
\sidebyside{figures/student_crop_lynae_before_input.png}{원본의 실제 입력}{figures/student_crop_lynae_after_input.png}{크롭 사본의 실제 입력}
첨부 예시처럼 윗부분을 덜 포함하고 아래쪽 얼굴과 옷이 함께 보이도록 잘랐다. 실제 중앙 입력에서 옷과 팔의 영역이 늘었으며, 풍선껌과 분홍/보라 계열 배경 일부는 그대로 남는다. 따라서 풍선껌이나 배경색 하나가 오분류의 원인이었다고 단정할 수 없다.
\end{minipage}\par\medskip

\subsection*{What this follow-up establishes and leaves open}
학생이 예상한 방향대로 두 크롭에서 모델의 선택이 정답으로 바뀌었다.
이는 이 두 사진의 구도를 바꾼 뒤 관찰한 결과이며, 각 캐릭터 전체로 일반화할 근거는 부족하다.
README는 각 변경을 여러 이미지에서 비교하라고 하며, 캐릭터마다 여러 장을 요구하지는 않는다.
이 두 사진과 아래 네 제공본은 학생이 선택한 크롭을 여러 이미지에서 비교한 기록이다.
학생은 아래 네 제공본 결과까지 확인한 뒤,
예상과 달리 점수가 낮아진 경우와 정사각형 입력에서 의상이 잘릴 가능성에 대한 해석을 직접 제시했다.
그 원문과 실제 입력의 관찰은 아래 최신 학생 해석 절에 기록했다.
배경, 캐릭터의 크기/위치와 보이는 얼굴/의상이 함께 바뀌었으므로 하나의 특징만을 분리한 결론을 내리지 않는다.

\clearpage
\subsection*{Measured output and demo comparison}
다음은 정확한 크롭 좌표, 전후 답, 전체 범주 점수와 파일 보존 확인을 코드가 출력한 원문이다.
재현 코드와 학생 요청/좌표는 ZIP의 \path{outputs/student_crops_code.py}와
\path{outputs/student_crops_request.json}에도 보존했다.
\outputfile{outputs/student_crops_output.txt}
로컬 데모에 두 원본과 두 크롭본을 실제로 업로드하고,
현재 내보낸 같은 모델과 실제 입력, 답과 점수를 확인한 기록은 다음과 같다.
\outputfile{outputs/student_crops_browser_summary.txt}
\par\noindent\begin{minipage}{\linewidth}
\sidebyside{figures/student_crop_browser_camellya_IMG_2554.png}{카멜리아 크롭본 업로드}{figures/student_crop_browser_lynae_IMG_0860.png}{린네 크롭본 업로드}
\end{minipage}\par\medskip
전체 브라우저 측정 원문과 배열은 ZIP의 \path{outputs/student_crops_browser_output.txt}와
\path{outputs/student_crops_browser_results.json}에 보존했다.



\clearpage
\subsection*{Student-provided follow-up crop files}
\label{p7:student_provided_crops}
앞서 정답으로 분류된 사진 중 정답 점수가 낮았던 사진을 대상으로,
학생이 직접 만든 서로 다른 구도의 크롭본을 다음 순서로 전달했다.
\begin{quote}
``순서대로 iimg\_0845, img\_0806, img\_0818, img\_0844 야''
\end{quote}

정확한 파일을 다시 제공하면서 순서를 다음과 같이 확인했다.
\begin{quote}
``각각 순서대로야''
\end{quote}

\texttt{iimg\_0845}는 기존 \texttt{IMG\_0845.jpg}에 대응시켰다.
이번에는 네 제공 파일 모두에 접근할 수 있어 위 순서대로 측정했다.
앞서 \texttt{IMG\_0844} 한 장만 측정하고 세 파일을 대기한 기록은
\path{outputs/supplied_crops_initial.json} 등 \texttt{\_initial} 파일에 과거 증거로 보존했다.
제공 파일을 바이트 그대로 복사했고, 새 크롭 좌표를 추정하거나 추가 편집하지 않았다.
같은 \texttt{expanded} ONNX 모델과 기존 RGB/resize/중앙 크롭 전처리를 사용해
각 원본과 제공본을 한 번씩 평가했다.

\begin{center}\begin{tabular}{lrr}
\toprule 사진 & 원본의 정답 점수 & 제공본의 정답 점수 \\
\midrule
린네 \texttt{IMG\_0845} & 65.7\% & 87.5\% \\
에이메스 \texttt{IMG\_0806} & 69.6\% & 100.0\% \\
에이메스 \texttt{IMG\_0818} & 75.8\% & 100.0\% \\
린네 \texttt{IMG\_0844} & 80.9\% & 68.5\% \\
\bottomrule\end{tabular}\end{center}
정답 라벨의 softmax 점수는 3개에서 높아지고 1개에서 낮아졌다.
선택된 원본은 4/4, 제공본은 4/4를 맞혔다.
모델이 실제로 선택한 이름과 그 이름의 점수는 각 그림에 표시했다.
\texttt{100.0\%}는 한 자리로 반올림한 softmax 점수이며 완전한 확신이나 항상 맞는다는 뜻이 아니다.
이는 본래 정답인 사진 중에서 선택한 사례의 전후 비교이며 독립된 새 평가 정확도가 아니다.
개인 평가 원본의 기존 결과는 바꾸지 않았다.
제공본도 여러 사진의 크롭 전후 비교이며, 사진별 구도와 프레이밍이 달라 동일한 면적이나 하나의 특징만을 통제한 비교는 아니다.
보이는 얼굴/의상, 대상의 크기와 위치, 배경, 스크린샷 저장 과정이 함께 달라질 수 있어
원본의 정확한 RGB 픽셀 부분영역이나 단일 특징의 인과 효과라고 주장하지 않는다.
이전 카멜리아/린네 두 오답의 학생 지정 크롭과 이번 네 제공본은 별도 기록이다.
학생은 점수가 오른 결과는 예상에 맞았으나 하락은 예상하지 못했다고 설명했다.
그리고 세로 제공본의 중앙 정사각형 입력에서 의상이 잘릴 수도 있다는 잠정 추측을 제시했다.
아래에 그 원문을 보존하며, 새로운 변환 실행을 요청한 것으로 해석하지 않는다.
최초 두 사례와 네 제공본을 합해 학생이 선택한 여섯 사진에서 크롭 전후를 비교했다.
캐릭터별 사진 수가 다르다는 이유만으로 여러 이미지 비교가 미완료라고 판단하지 않는다.
이 결과와 학생의 해석은 여러 사진의 구도 변경에 대한 증거이며, 특정 특징 하나의 인과 효과를 확정하지는 않는다.

\clearpage
\subsection*{Provided framing: \texttt{IMG\_0845.jpg}}
\par\noindent\begin{minipage}{\linewidth}
\sidebyside{figures/supplied_crop_IMG_0845_original.jpg}{원본: 린네 65.7\%}{figures/supplied_crop_IMG_0845_changed.png}{제공본: 린네 87.5\%}
위는 전체 원본과 학생 제공본, 아래는 각각 모델에 실제로 들어간 224 정사각형 입력이다.
\sidebyside{figures/supplied_crop_IMG_0845_before_input.png}{원본의 실제 224 입력}{figures/supplied_crop_IMG_0845_after_input.png}{제공본의 실제 224 입력}
\end{minipage}\par\medskip

\clearpage
\subsection*{Provided framing: \texttt{IMG\_0806.jpg}}
\par\noindent\begin{minipage}{\linewidth}
\sidebyside{figures/supplied_crop_IMG_0806_original.jpg}{원본: 에이메스 69.6\%}{figures/supplied_crop_IMG_0806_changed.png}{제공본: 에이메스 100.0\%}
위는 전체 원본과 학생 제공본, 아래는 각각 모델에 실제로 들어간 224 정사각형 입력이다.
\sidebyside{figures/supplied_crop_IMG_0806_before_input.png}{원본의 실제 224 입력}{figures/supplied_crop_IMG_0806_after_input.png}{제공본의 실제 224 입력}
\end{minipage}\par\medskip

\clearpage
\subsection*{Provided framing: \texttt{IMG\_0818.jpg}}
\par\noindent\begin{minipage}{\linewidth}
\sidebyside{figures/supplied_crop_IMG_0818_original.jpg}{원본: 에이메스 75.8\%}{figures/supplied_crop_IMG_0818_changed.png}{제공본: 에이메스 100.0\%}
위는 전체 원본과 학생 제공본, 아래는 각각 모델에 실제로 들어간 224 정사각형 입력이다.
\sidebyside{figures/supplied_crop_IMG_0818_before_input.png}{원본의 실제 224 입력}{figures/supplied_crop_IMG_0818_after_input.png}{제공본의 실제 224 입력}
\end{minipage}\par\medskip

\clearpage
\subsection*{Provided framing: \texttt{IMG\_0844.jpg}}
\par\noindent\begin{minipage}{\linewidth}
\sidebyside{figures/supplied_crop_IMG_0844_original.jpg}{원본: 린네 80.9\%}{figures/supplied_crop_IMG_0844_changed.png}{제공본: 린네 68.5\%}
위는 전체 원본과 학생 제공본, 아래는 각각 모델에 실제로 들어간 224 정사각형 입력이다.
\sidebyside{figures/supplied_crop_IMG_0844_before_input.png}{원본의 실제 224 입력}{figures/supplied_crop_IMG_0844_after_input.png}{제공본의 실제 224 입력}
\end{minipage}\par\medskip

\clearpage
\subsection*{Student interpretation after the crop comparisons}
\label{p7:student_crop_reflection}
학생은 두 오답의 지정 크롭과 네 제공본 비교 결과를 확인한 뒤 다음과 같이 설명했다.
\begin{quote}
``줄어들 걸 예상하지 못 했고 올라갈 건 다 예상에 맞았다. 린네인데 에이메스로 헷갈렸던 게 린네 인형탈의 옷이 잘 보이지 않았기에 옷까지 있는 전체를 크롭하면 올라갈 줄 알았다. 근데 다시 생각해보니 세로로 크롭했기 때문에 모델이 그걸 다시 정사각형으로 읽으면서 옷이 안 들어갔을 수도 있겠다는 생각을 했다.''
\end{quote}

이 설명의 에이메스 오분류 사례와 이번 점수 하락 사례는 서로 다른 사진이다.
앞서 \texttt{IMG\_0860.jpg}는 에이메스로 오분류되었고, 아래쪽을 포함한 지정 크롭 뒤 린네로 바뀌었다.
해당 정답 점수는 17.2\%에서 67.5\%로 올랐다.
반면 \texttt{IMG\_0844.jpg}는 원본과 제공본 모두 린네로 분류되었으며,
정답 점수가 80.9\%에서 68.5\%로 낮아진 경우다.

앞 페이지의 실제 224 입력을 보면 \texttt{IMG\_0844} 원본에는 아래쪽 흰 상의/소매와 팔이 비교적 넓게 보인다.
제공본의 실제 입력에서는 얼굴과 머리가 더 크게 차지하며, 흰 소매는 아래 가장자리에 좁게 남는다.
따라서 옷이 완전히 사라진 것은 아니지만 보이는 영역이 줄었다는 관찰은 학생의 잠정 추측과 관련된다.
별도 \texttt{IMG\_0860}의 지정 크롭에서는 실제 입력에 상의와 몸통 일부가 더 많이 들어온다.

학생의 최신 추측은 전체 제공 파일에 옷이 있어도 중앙 정사각형 입력에서는 일부가 잘릴 수 있다는 것이다.
머리와 옷이 단서일 것이라는 기존 예상에, 모델이 실제로 보는 영역을 확인할 필요가 있다는 해석을 더했다.
학생은 ``안 들어갔을 수도''라고 표현했으며, 이는 원인을 확정한 결론이 아니다.
의상, 얼굴/머리의 상대적 크기, 배경과 위치가 함께 바뀌었으므로 의상 포함량 하나의 효과를 입증하지 않는다.
학생이 별도의 확신 수치를 제시하거나 새 변환 실행을 지정한 것은 아니므로 이를 임의로 추가하지 않았다.
기존 요청 및 실제 파일 대응과 함께 이번 설명 원문은 \path{outputs/student_review_and_reflection.json}에도 보존했다.

\clearpage
\subsection*{Measured output from all received follow-up files}
다음은 네 제공본과 원본을 비교한 코드의 전체 출력이다. 대기 파일은 없다.
재현 코드, 학생 요청과 파일 대응, 파일 보존 확인 및 모든 범주 점수는 ZIP의
\path{outputs/supplied_crops_code.py}, \path{outputs/supplied_crops_request.json},
\path{outputs/supplied_crops.json}에 보존했다.
\outputfile{outputs/supplied_crops_output.txt}
같은 네 원본과 네 제공본을 실제 로컬 데모에 업로드해 입력 RGB와 답을 비교했다.
\outputfile{outputs/supplied_crops_browser_summary.txt}
전체 브라우저 측정 원문과 배열은 \path{outputs/supplied_crops_browser_output.txt}와
\path{outputs/supplied_crops_browser_results.json}에 보존했다.

\clearpage
\subsection*{Actual demo uploads of supplied framings}
\par\noindent\begin{minipage}{\linewidth}
\sidebyside{figures/supplied_crop_IMG_0845_browser_original.png}{\texttt{IMG\_0845} 원본 데모 업로드}{figures/supplied_crop_IMG_0845_browser_changed.png}{\texttt{IMG\_0845} 제공본 데모 업로드}
\end{minipage}\par\medskip
\par\noindent\begin{minipage}{\linewidth}
\sidebyside{figures/supplied_crop_IMG_0806_browser_original.png}{\texttt{IMG\_0806} 원본 데모 업로드}{figures/supplied_crop_IMG_0806_browser_changed.png}{\texttt{IMG\_0806} 제공본 데모 업로드}
\end{minipage}\par\medskip

\clearpage
\subsection*{Actual demo uploads of supplied framings}
\par\noindent\begin{minipage}{\linewidth}
\sidebyside{figures/supplied_crop_IMG_0818_browser_original.png}{\texttt{IMG\_0818} 원본 데모 업로드}{figures/supplied_crop_IMG_0818_browser_changed.png}{\texttt{IMG\_0818} 제공본 데모 업로드}
\end{minipage}\par\medskip
\par\noindent\begin{minipage}{\linewidth}
\sidebyside{figures/supplied_crop_IMG_0844_browser_original.png}{\texttt{IMG\_0844} 원본 데모 업로드}{figures/supplied_crop_IMG_0844_browser_changed.png}{\texttt{IMG\_0844} 제공본 데모 업로드}
\end{minipage}\par\medskip

\clearpage
\subsection*{Student-selected hair masking on all personal images}
\label{sec:student-hair}
이 실험은 학생이 직접 고른 머리카락 가림을 개인 평가 사진 전체에 적용한 후속 탐구다.
앞선 크롭은 구도와 대상의 크기, 배경을 함께 바꾸었으므로,
이번에는 같은 모델 입력에서 지정한 머리카락 픽셀만 회색으로 바꾸었다.
모델을 다시 학습하거나 바꾸지 않았고, 가림 이후 답을 보고 사진을 골라내거나 좌표를 수정하지 않았다.

\subsubsection*{Student hypothesis, scope and execution requests}
학생이 선택한 가설과 변경, 대상 범위와 실행 방법을 다음 순서의 원문으로 보존한다.

\begin{quote}
``음 이미 말했지 않나 머리스타일을 보고 구별할 것 같고, 사진에서 머리카락만 제외해서 재확인하겠다''
\end{quote}

\begin{quote}
``24장 전부. 응 머리카락 부분만 회색으로 가리는 방식''
\end{quote}

\begin{quote}
``Python으로 정확한 픽셀 가림''
\end{quote}

\subsubsection*{Controlled pixels and the actual model view}
각 원본을 RGB로 읽고 짧은 변을 256으로 bilinear 리사이즈한 뒤 중앙 256 정사각형 사본을 만들었다.
이 사본을 기존 \texttt{prepare\_image}로 처리한 배열은 전체 원본을 처리한 224 입력 배열과 정확히 같았다.
이 전처리 동등성을 먼저 확인한 사본에서만 머리카락 영역을 표시했다.
따라서 아래 ``원본 256''은 전체 촬영 파일이 아니라, 기존 원본의 모델 입력을 그대로 재현하는 표시용 사본이다.
``실제 입력 224''는 그 사본의 중앙에서 모델이 보는 영역이며, 두 단계의 이미지를 함께 보여 준다.

에이전트가 답을 계산하기 전에 보이는 대상 캐릭터의 머리카락을 수동 다각형으로 표시하고 시각 검수했다.
눈, 얼굴, 손, 옷, 장식 및 다른 캐릭터는 제외 영역으로 보호했다.
마스크 안은 RGB(128, 128, 128)로 채웠으며, 밖은 원래 픽셀과 정확히 같음을 검사했다.
원본 개인 파일, 학습 자료, 기존 크롭 실험과 내보낸 모델은 바꾸지 않았다.
재현 명령은 다음과 같다.
\begin{verbatim}
uv run --frozen python -B src/hair_mask_experiment.py
  --config config/student_hair_masks.json
\end{verbatim}
위 두 줄은 한 명령으로 이어서 실행한다.
고정 좌표와 검수 기록은 \path{outputs/student_hair_masks_config.json},
실행 코드는 \path{outputs/student_hair_masks_code.py}에 보존했다.
모든 범주의 점수와 마스크 픽셀 수는 \path{outputs/student_hair_masks.json} 및 뒤의 전체 출력에 있다.

흐릿한 경계와 아주 작은 인쇄 그림은 완전하게 분할할 수 없었다.
예를 들어 린네 \texttt{IMG\_0871}의 자동차 보닛 그림, 카멜리아 \texttt{IMG\_0583}의 장식 사진,
다른 사진의 가는 머리카락 끝과 장식 사이 경계 일부는 남겼다.
회색으로 채워도 바깥 윤곽, 위치와 머리의 크기는 남는다.
이 조작을 ``머리스타일을 완전히 제거한 실험''으로 해석하지 않는다.
회색 영역 자체도 평소 사진에 없는 자극이므로 결과에는 가림의 부작용이 함께 포함될 수 있다.

\subsubsection*{Measured result}
\begin{center}\small
\begin{tabular}{lrrr}
\toprule
범주 & 원본 정답 & 머리 가림 후 정답 & 예측 범주 변경 \\
\midrule

린네 & 5/6 (83.3\%) & 4/6 (66.7\%) & 1 \\

양양 & 6/6 (100.0\%) & 5/6 (83.3\%) & 1 \\

에이메스 & 6/6 (100.0\%) & 6/6 (100.0\%) & 0 \\

카멜리아 & 5/6 (83.3\%) & 5/6 (83.3\%) & 0 \\

전체 & 22/24 (91.7\%) & 20/24 (83.3\%) & 2 \\

\bottomrule
\end{tabular}
\end{center}
원본 개인 평가의 정답은 그대로 22/24장(91.7\%)이며, 가림 후에는 20/24장(83.3\%)이었다.
양양 \texttt{IMG\_0810}은 린네로, 린네 \texttt{IMG\_0850}은 에이메스로 바뀌었다.
기존 두 오답은 가림 후에도 오답이었고, 새로운 두 오답이 추가되었다.
실제 캐릭터에 준 점수는 19장에서 하락하고 5장에서 상승했다.
증감은 반올림 전 확률로 집계했으므로, 표에 같은 값으로 인쇄되어도 작은 증감이 있을 수 있다.
정답 점수는 해당 사진 한 장에서 실제 범주에 준 확률이며, 여러 사진의 정답률과 구분한다.

\subsubsection*{Student interpretation after seeing the result}
학생은 출력 결과를 확인한 뒤 다음과 같이 설명했다. 실험 전 예상으로 소급하지 않는다.

\begin{quote}
``카멜리아 같은 경우 예시 사진의 옷이 다른 사진이 없다. 그래서 머리카락을 가려도 옷을 보고 카멜리아라고 유추할 수 있던 것 같다. 나머지는 예시 사진이 조금씩 다르게 생겼기 때문에 머리카락이 가려지면 큰 영향을 받은 것 같다.''
\end{quote}

카멜리아의 예측이 유지된 결과를, 학생은 옷이 남아 있다는 가능한 설명과 연결했다.
양양과 린네의 일부 사진은 예측이 달라졌고 에이메스도 정답 점수가 낮아졌으므로,
머리카락의 내부 색과 무늬가 일부 판단에 영향을 주었다는 해석과 부합한다.
다만 옷만으로 구별했는지, 머리 윤곽이나 장식·얼굴·배경을 함께 사용했는지는 이번 가림 하나로 분리할 수 없다.
카멜리아의 유지와 다른 범주의 변화가 의상 차이 때문에 발생했다는 설명도 학생의 잠정 해석이다.
학생이 옷을 별도로 바꾸는 새 실험을 선택한 것은 아니므로 추가 조작을 임의로 수행하지 않았다.
학생 원문과 수신 시점은 \path{outputs/student_hair_masks_reflection.json}에 보존했다.

\clearpage
\subsubsection*{All cases: predicted labels and true-class scores}
예측 답 열은 범주와 그 범주의 점수이며, 정답 점수 열은 실제 범주의 점수다.
\begin{center}\scriptsize
\setlength{\tabcolsep}{3pt}
\begin{tabular}{llllrr}
\toprule
실제 범주 & 파일 & 원본 예측 답 & 가림 후 예측 답 & 원본 정답 점수 & 후 정답 점수 \\
\midrule

양양 & \texttt{IMG\_0804.jpg} & 양양 98.6\% & 양양 97.3\% & 98.6\% & 97.3\% \\

양양 & \texttt{IMG\_0809.jpg} & 양양 98.2\% & 양양 99.5\% & 98.2\% & 99.5\% \\

양양 & \texttt{IMG\_0810.jpg} & 양양 93.7\% & 린네 78.9\% & 93.7\% & 20.2\% \\

양양 & \texttt{IMG\_0812.jpg} & 양양 96.6\% & 양양 89.3\% & 96.6\% & 89.3\% \\

양양 & \texttt{IMG\_1390.JPG} & 양양 99.5\% & 양양 88.2\% & 99.5\% & 88.2\% \\

양양 & \texttt{IMG\_1392.JPG} & 양양 99.8\% & 양양 83.0\% & 99.8\% & 83.0\% \\

\midrule

카멜리아 & \texttt{IMG\_0583.jpg} & 카멜리아 99.8\% & 카멜리아 99.6\% & 99.8\% & 99.6\% \\

카멜리아 & \texttt{IMG\_0584.jpg} & 카멜리아 100.0\% & 카멜리아 100.0\% & 100.0\% & 100.0\% \\

카멜리아 & \texttt{IMG\_0834.jpg} & 카멜리아 94.9\% & 카멜리아 97.4\% & 94.9\% & 97.4\% \\

카멜리아 & \texttt{IMG\_0867.jpg} & 카멜리아 99.5\% & 카멜리아 99.5\% & 99.5\% & 99.5\% \\

카멜리아 & \texttt{IMG\_2554.jpg} & 에이메스 80.8\% & 에이메스 85.2\% & 0.9\% & 0.7\% \\

카멜리아 & \texttt{IMG\_2555.jpg} & 카멜리아 100.0\% & 카멜리아 100.0\% & 100.0\% & 100.0\% \\

\midrule

린네 & \texttt{IMG\_0844.jpg} & 린네 80.9\% & 린네 53.5\% & 80.9\% & 53.5\% \\

린네 & \texttt{IMG\_0845.jpg} & 린네 65.7\% & 린네 37.1\% & 65.7\% & 37.1\% \\

린네 & \texttt{IMG\_0850.jpg} & 린네 87.5\% & 에이메스 85.9\% & 87.5\% & 10.8\% \\

린네 & \texttt{IMG\_0860.jpg} & 에이메스 45.6\% & 에이메스 50.6\% & 17.2\% & 4.4\% \\

린네 & \texttt{IMG\_0871.jpg} & 린네 89.2\% & 린네 78.5\% & 89.2\% & 78.5\% \\

린네 & \texttt{IMG\_1417.JPG} & 린네 96.6\% & 린네 92.8\% & 96.6\% & 92.8\% \\

\midrule

에이메스 & \texttt{IMG\_0215.PNG} & 에이메스 100.0\% & 에이메스 50.6\% & 100.0\% & 50.6\% \\

에이메스 & \texttt{IMG\_0553.jpg} & 에이메스 99.9\% & 에이메스 84.7\% & 99.9\% & 84.7\% \\

에이메스 & \texttt{IMG\_0802.jpg} & 에이메스 98.8\% & 에이메스 98.7\% & 98.8\% & 98.7\% \\

에이메스 & \texttt{IMG\_0806.jpg} & 에이메스 69.6\% & 에이메스 54.2\% & 69.6\% & 54.2\% \\

에이메스 & \texttt{IMG\_0818.jpg} & 에이메스 75.8\% & 에이메스 61.9\% & 75.8\% & 61.9\% \\

에이메스 & \texttt{IMG\_0869.jpg} & 에이메스 81.1\% & 에이메스 71.1\% & 81.1\% & 71.1\% \\


\bottomrule
\end{tabular}
\end{center}

\clearpage
\subsubsection*{양양 / 비교 1}
왼쪽 두 열은 모델 입력과 동등한 256 사본의 원본/가림본이고, 오른쪽 두 열은 각각의 실제 중앙 224 입력이다.
각 사진 아래에 해당 입력에 대한 모델의 답과 실제 범주의 점수를 적었다.
\par\noindent \picturefile{figures/hair_yangyang_1.png}\par

\clearpage
\subsubsection*{양양 / 비교 2}
왼쪽 두 열은 모델 입력과 동등한 256 사본의 원본/가림본이고, 오른쪽 두 열은 각각의 실제 중앙 224 입력이다.
각 사진 아래에 해당 입력에 대한 모델의 답과 실제 범주의 점수를 적었다.
\par\noindent \picturefile{figures/hair_yangyang_2.png}\par

\clearpage
\subsubsection*{카멜리아 / 비교 1}
왼쪽 두 열은 모델 입력과 동등한 256 사본의 원본/가림본이고, 오른쪽 두 열은 각각의 실제 중앙 224 입력이다.
각 사진 아래에 해당 입력에 대한 모델의 답과 실제 범주의 점수를 적었다.
\par\noindent \picturefile{figures/hair_camellya_1.png}\par

\clearpage
\subsubsection*{카멜리아 / 비교 2}
왼쪽 두 열은 모델 입력과 동등한 256 사본의 원본/가림본이고, 오른쪽 두 열은 각각의 실제 중앙 224 입력이다.
각 사진 아래에 해당 입력에 대한 모델의 답과 실제 범주의 점수를 적었다.
\par\noindent \picturefile{figures/hair_camellya_2.png}\par

\clearpage
\subsubsection*{린네 / 비교 1}
왼쪽 두 열은 모델 입력과 동등한 256 사본의 원본/가림본이고, 오른쪽 두 열은 각각의 실제 중앙 224 입력이다.
각 사진 아래에 해당 입력에 대한 모델의 답과 실제 범주의 점수를 적었다.
\par\noindent \picturefile{figures/hair_lynae_1.png}\par

\clearpage
\subsubsection*{린네 / 비교 2}
왼쪽 두 열은 모델 입력과 동등한 256 사본의 원본/가림본이고, 오른쪽 두 열은 각각의 실제 중앙 224 입력이다.
각 사진 아래에 해당 입력에 대한 모델의 답과 실제 범주의 점수를 적었다.
\par\noindent \picturefile{figures/hair_lynae_2.png}\par

\clearpage
\subsubsection*{에이메스 / 비교 1}
왼쪽 두 열은 모델 입력과 동등한 256 사본의 원본/가림본이고, 오른쪽 두 열은 각각의 실제 중앙 224 입력이다.
각 사진 아래에 해당 입력에 대한 모델의 답과 실제 범주의 점수를 적었다.
\par\noindent \picturefile{figures/hair_aemeath_1.png}\par

\clearpage
\subsubsection*{에이메스 / 비교 2}
왼쪽 두 열은 모델 입력과 동등한 256 사본의 원본/가림본이고, 오른쪽 두 열은 각각의 실제 중앙 224 입력이다.
각 사진 아래에 해당 입력에 대한 모델의 답과 실제 범주의 점수를 적었다.
\par\noindent \picturefile{figures/hair_aemeath_2.png}\par

\clearpage
\subsubsection*{Complete printed hair-mask measurement and preservation checks}
아래는 고정된 기존 ONNX 모델로 원본과 가림본을 한 번씩 평가한 전체 코드 출력이다.
별도 학습, 새 모델 선택이나 점수가 좋은 가림본만 골라 보고하는 과정은 없다.
이번 새 가림본은 Python에서 내보낸 ONNX 모델로 평가했다.
브라우저 제어 연결이 제공되지 않아 새 원본/가림본의 실제 데모 업로드 비교는 이번 실행에서 확인하지 못했다.
앞선 크롭의 브라우저 업로드 기록과 구분하며, README Problem 7의 데모 비교 단계는 이 새 변환에 대해서는 남아 있다.
\outputfile{outputs/student_hair_masks_output.txt}
기존 개인 평가 기록과 모든 원본 답의 대조, 파일 보존 검사는 다음과 같다.
\outputfile{outputs/student_hair_masks_verification.txt}

\section*{Appendix. Additional measured evidence}
현재 \texttt{expanded}의 원본·변환 예측과 합계, 전처리/분할/기존 파일 보존 검사는 다음 출력과 같다.
\outputfile{outputs/expanded_analysis_output.txt}
\outputfile{outputs/expanded_dataset_audit.txt}

아래는 과거 \texttt{final}의 분석과 분할 검사 및 기존 테스트 결과이며 그대로 보존했다.
\outputfile{outputs/analysis_output.txt}
\outputfile{outputs/split_content_audit.txt}
\outputfile{outputs/tests.txt}

\section*{Appendix. Historical final demo}
아래 출력과 그림은 이번 확장 전에 사용한 \texttt{final} 모델의 기록이다.
\outputfile{outputs/export.txt}
\outputfile{outputs/browser_demo.txt}
\outputfile{outputs/browser_ablations.txt}
\par\noindent \picturefile{figures/browser_selftest.png}\par
\par\noindent \picturefile{figures/browser_ood.png}\par
\section*{Appendix. Complete output of check.py: current expanded}
아래는 현재 내보낸 \texttt{expanded} 모델에 대한 공식 검사 전체 출력이다.
학생이 커밋을 직접 하겠다고 요청했으므로 커밋 여부 경고를 성능 오류로 해석하지 않는다.
\outputfile{outputs/expanded_check.txt}

\section*{Appendix. Complete output of check.py: historical final}
아래는 확장 전 \texttt{final} 모델의 공식 검사 전체 출력으로 보존한 과거 기록이다.
\outputfile{outputs/check.txt}
\end{document}
